%==================================================================
%  PROJECT KUMQUAT  --  THESIS_TEMPLATE.tex
%  The Engineer's Blueprint for Veterinary Medicine & Stray Triage
%  "Vet-level grounding for non-vets."
%------------------------------------------------------------------
%  COMPILE WITH : xelatex THESIS_TEMPLATE.tex   (run TWICE for TOC)
%                 (lualatex also works)
%  REQUIRES     : A Unicode TeX engine (XeLaTeX / LuaLaTeX) and an
%                 Arabic-capable OpenType font. Default = Amiri (free).
%                 Free alternatives: "Scheherazade New",
%                 "Noto Naskh Arabic". Change the \newfontfamily line
%                 below if Amiri is not installed.
%------------------------------------------------------------------
%  NOTE ON ARABIC RENDERING
%  The project specification references arabtex. Under a modern Unicode
%  engine the robust, equivalent system is polyglossia + fontspec, which
%  renders native inline Arabic flawlessly via \textarabic{...} (wrapped
%  here in the convenience macro \ar{...}). A legacy pdfLaTeX + arabtex
%  fallback is documented in the comment block at the very bottom of
%  this file.
%==================================================================
\documentclass[12pt,oneside]{book}

%------------------------- Page geometry --------------------------
\usepackage[a4paper,top=1in,bottom=1in,left=1.25in,right=1in,headheight=15pt]{geometry}

%------------------------- Mathematics ----------------------------
\usepackage{amsmath}
\usepackage{amssymb}

%------------------------- Graphics & color -----------------------
\usepackage{xcolor}
\usepackage{graphicx}
\usepackage{textcomp}        % \textdegree and friends
\usepackage{booktabs}        % professional tables
\usepackage{longtable}
\usepackage{array}
\usepackage{enumitem}        % tunable lists
\usepackage{tcolorbox}       % callout boxes (engineering "panels")

\definecolor{kumdeep}{HTML}{0B3D5C}    % deep system blue
\definecolor{kumaccent}{HTML}{E08A2B}  % kumquat orange
\definecolor{kumgrey}{HTML}{5B6770}
\definecolor{kumpanel}{HTML}{F4F7F9}

%------------------------- Headings (titlesec) --------------------
\usepackage{titlesec}
\titleformat{\chapter}[display]
  {\normalfont\huge\bfseries\color{kumdeep}}
  {\filright\large\sffamily\color{kumaccent}\MakeUppercase{Chapter \thechapter}}
  {1ex}
  {\titlerule[1.2pt]\vspace{1ex}\filright}
\titlespacing*{\chapter}{0pt}{10pt}{26pt}

\titleformat{\section}
  {\normalfont\Large\bfseries\color{kumdeep}}{\thesection}{0.75em}{}
\titleformat{\subsection}
  {\normalfont\large\bfseries\color{kumgrey}}{\thesubsection}{0.75em}{}

%------------------------- Headers / footers ----------------------
\usepackage{fancyhdr}
\pagestyle{fancy}
\fancyhf{}
\renewcommand{\headrulewidth}{0.4pt}
\fancyhead[L]{\small\sffamily\color{kumgrey}Project Kumquat}
\fancyhead[R]{\small\sffamily\color{kumgrey}\leftmark}
\fancyfoot[C]{\thepage}

%------------------------- Hyperlinks -----------------------------
%  (Load hyperref BEFORE polyglossia so bidi initialises last.)
\usepackage{hyperref}
\hypersetup{
  colorlinks=true,
  linkcolor=kumdeep,
  urlcolor=kumaccent,
  citecolor=kumdeep,
  pdftitle={The Engineer's Blueprint for Veterinary Medicine and Stray Triage},
  pdfauthor={Project Kumquat Contributors},
  pdfsubject={Veterinary systems architecture and field triage},
  pdfkeywords={veterinary, triage, panleukopenia, systems engineering, rescue}
}

%------------------------- Multilingual ---------------------------
%  polyglossia pulls in bidi for Arabic; keep it LAST among packages.
\usepackage{fontspec}
\usepackage{polyglossia}
\setmainlanguage{english}
\setotherlanguage{arabic}
%  Swap "Amiri" below for any installed Arabic OpenType font.
\newfontfamily\arabicfont[Script=Arabic]{Amiri}
%  Convenience macro: inline Arabic inside English text.
\newcommand{\ar}[1]{\textarabic{#1}}

%------------------------- Reusable panel -------------------------
\newtcolorbox{kumbox}[1]{
  colback=kumpanel, colframe=kumdeep, coltitle=white,
  fonttitle=\bfseries\sffamily, title={#1},
  boxrule=0.6pt, arc=2pt, left=6pt, right=6pt, top=4pt, bottom=4pt}

%==================================================================
\begin{document}
\frontmatter

%------------------------- Title page -----------------------------
\begin{titlepage}
\centering
\vspace*{1.2cm}
{\sffamily\color{kumaccent}\large PROJECT KUMQUAT \quad$\cdot$\quad CODENAME KUMQUAT\par}
\vspace{0.4cm}
\rule{\linewidth}{1.2pt}\par
\vspace{0.8cm}
{\Huge\bfseries\color{kumdeep} The Engineer's Blueprint for\\[0.3cm] Veterinary Medicine \& Stray Triage\par}
\vspace{0.6cm}
{\Large\itshape A System-Architecture Approach to Biology\par}
\vspace{0.3cm}
{\large \textit{Vet-level grounding for non-vets}\par}
\vspace{0.8cm}
\rule{\linewidth}{0.5pt}\par
\vspace{1.4cm}
{\large A structured curriculum for field rescuers, wildlife emergency\\
responders, and bio-engineers who require deterministic triage\\
frameworks under pressure.\par}
\vspace{1.6cm}
{\large Feline Panleukopenia \quad \ar{طاعون القطط الفيروسي}\par}
\vfill
{\sffamily\color{kumgrey} In memory of Kumquat (Koki) --- the baseline that taught us the cost of latency.\par}
\vspace{0.6cm}
{\large 2026 \quad$\cdot$\quad CC-BY-SA-4.0\par}
\end{titlepage}

%------------------------- Disclaimer -----------------------------
\chapter*{Disclaimer \& Scope}
\addcontentsline{toc}{chapter}{Disclaimer \& Scope}
\thispagestyle{fancy}

\begin{kumbox}{GROUNDING GUARDRAIL}
This document does \textbf{not} replace a licensed Doctor of Veterinary
Medicine (DVM) degree, nor is it a casual pet-owner guide. It provides
\emph{vet-level grounding for non-vets}. Every diagnostic path described
here terminates at the same instruction: \textbf{stabilize, then hand off
to a professional veterinarian}. Nothing in this work authorizes
diagnosis, prescription, or surgery.
\end{kumbox}

\noindent This curriculum is engineered for individuals who require
high-fidelity, academically grounded veterinary knowledge --- such as
field rescuers, wildlife emergency responders, and bio-engineers ---
without navigating a full clinical licensing array. It rejects dense,
traditional medical literature in favour of a system-architecture model:
the animal body is treated as an integrated hardware stack, organs are
specialized sub-routines, the vascular network is fluid plumbing, and the
immune system is a dynamic, multi-tier firewall. When a system goes down,
we do not guess --- we debug.

%------------------------- Acknowledgements -----------------------
\chapter*{Acknowledgements}
\addcontentsline{toc}{chapter}{Acknowledgements}
\thispagestyle{fancy}

Project Kumquat was born out of a critical baseline breakdown. When a
tiny, vulnerable ginger kitten was rescued from an industrial zone, we
encountered the unforgiving, hardcoded realities of nature: a fast-moving,
hostile pathogen known as Feline Panleukopenia
(\ar{طاعون القطط الفيروسي}). In the field, time translates directly to
survival. Traditional veterinary resources are locked behind dense,
academic text walls that fail to serve rescuers, software engineers, and
logical problem-solvers who require deterministic triage frameworks under
pressure. This project honors the memory of that tiny fighter, Kumquat
(affectionately known as Koki). It rejects unstructured medical literature
in favor of a System Architecture Approach to Biology. We treat the animal
body as an integrated hardware stack: organs are specialized sub-routines,
the vascular network is fluid plumbing, and the immune system is a
dynamic, multi-tier fire-wall. When a system goes down, we do not
guess---we debug.

%------------------------- Abstract -------------------------------
\chapter*{Abstract}
\addcontentsline{toc}{chapter}{Abstract}
\thispagestyle{fancy}

This thesis presents a three-phase curriculum that re-frames foundational
veterinary science, pathophysiology, and emergency triage through the
lens of systems engineering. \textbf{Phase 1 (Foundational Systems
Architecture)} establishes healthy operating parameters and the
tissue-level production plants of the body --- bone marrow
(\ar{نخاع العظم}), the lymphatic transit network
(\ar{الجهاز اللمفاوي}), and localized filter checkpoints, the lymph nodes
(\ar{عقدة ليمفاوية}). \textbf{Phase 2 (Pathophysiology \& System Fault
Trees)} formalizes how viral, bacterial, and toxic payloads execute
code-injection attacks that drive cellular degradation, immune wipe
(Feline Panleukopenia, \ar{طاعون القطط الفيروسي}), gut-barrier breach,
and secondary systemic sepsis with liver shutdown. \textbf{Phase 3 (Field
Triage \& Runtime Debugging)} specifies a deterministic, step-by-step
diagnostic loop covering airway and breathing, thermal defense, capillary
refill time via the gum barrier, and fluid-collapse management prior to
professional veterinary handoff. The combined study load is approximately
114 hours. The document is bilingual: English is primary, with exact
Arabic scientific terminology embedded inline throughout.

\tableofcontents

%==================================================================
\mainmatter

%------------------------- Chapter 1: Intro -----------------------
\chapter{Introduction}

\section{The Problem of Latency}
In a field rescue, the gap between a correct decision and a guess is
measured in minutes, and minutes are measured in lives. Conventional
veterinary literature is optimized for the classroom and the clinic, not
for a responder kneeling over a collapsing animal in an industrial lot.
The information exists, but it is locked behind dense prose and assumes a
reader who already holds a clinical license. Project Kumquat exists to
close that gap with a deterministic, engineering-native framework.

\section{The Core Abstraction: Biology as a Hardware Stack}
We hold one mapping consistent across every chapter. Learn it once; reuse
it everywhere.

\begin{center}
\renewcommand{\arraystretch}{1.35}
\begin{tabular}{>{\bfseries}l l l}
\toprule
\textbf{Biological System} & \textbf{Engineering Analogue} & \textbf{Arabic} \\
\midrule
Bone marrow      & Cellular production plant (WBC factory) & \ar{نخاع العظم} \\
Lymphatic system & Transit network / message bus           & \ar{الجهاز اللمفاوي} \\
Lymph node       & Filter checkpoint / firewall node        & \ar{عقدة ليمفاوية} \\
Immune system    & Dynamic, multi-tier firewall             & \ar{الجهاز المناعي} \\
Vascular network & Fluid plumbing / pressurized pipework     & \ar{الجهاز الدوراني} \\
Liver            & Detox \& waste-processing daemon          & \ar{الكبد} \\
Gut barrier      & Perimeter wall / I/O membrane             & \ar{حاجز الأمعاء} \\
\bottomrule
\end{tabular}
\end{center}

\section{Method: We Debug, We Do Not Guess}
Each phase of this curriculum maps onto a recognizable engineering
discipline. Phase 1 is the \emph{specification sheet}: the documented,
healthy operating parameters of the system. Phase 2 is \emph{root-cause
analysis}: reading crash logs and constructing fault trees that trace a
failure from ingress to systemic collapse. Phase 3 is \emph{runtime
debugging}: a live, deterministic control loop executed under pressure,
whose single exit condition is a clean handoff to a professional
veterinarian.

\section{How to Read This Document}
Chapters proceed strictly in phase order. Do not skip the specification
sheet --- a fault cannot be detected without a known-good baseline.
Bilingual terminology is embedded inline so that a reader can cross
reference English and Arabic clinical vocabulary without leaving the page.

%------------------------- Chapter 2: Phase 1 ---------------------
\chapter{Phase 1 --- Foundational Systems Architecture}

\begin{kumbox}{PHASE 1 OBJECTIVE \quad ($\approx$ 36 study hours)}
Establish the system specification: healthy operating parameters and the
tissue-level production plants that manufacture and filter the body's
cellular workforce. \textbf{Exit criterion:} recite normal baselines from
memory and locate where each blood-cell line is produced and filtered.
\end{kumbox}

\section{The Hardware Stack Overview}
Before any fault can be diagnosed, the responder must hold a block diagram
of the healthy system: the major organ sub-routines, the pressurized
plumbing that connects them, and the firewall that defends them. This
module builds that top-level map and fixes the vocabulary used for the
remainder of the curriculum.

\section{Cellular Production Plants: Bone Marrow \texorpdfstring{(\ar{نخاع العظم})}{}}
Bone marrow (\ar{نخاع العظم}) is the body's primary manufacturing plant.
It produces the three critical cell lines:
\begin{itemize}[leftmargin=2em]
  \item \textbf{Red cells (erythrocytes)} --- the oxygen freight network.
  \item \textbf{White cells (leukocytes)} --- the firewall's active
        defenders. Their production rate is the single most important
        number in Phase 2.
  \item \textbf{Platelets (thrombocytes)} --- the automated leak-sealing
        subsystem.
\end{itemize}
Because the marrow is a high-mitosis core (cells divide rapidly), it is
also a \emph{high-value target} for any pathogen that hijacks cell
division --- a vulnerability formalized in Chapter 3.

\section{The Lymphatic Transit Network \texorpdfstring{(\ar{الجهاز اللمفاوي})}{}}
The lymphatic system (\ar{الجهاز اللمفاوي}) is the body's secondary
transit network: a one-way message bus that collects interstitial fluid,
transports immune cells, and routes antigens to checkpoints for
inspection. Unlike the pressurized vascular loop, it relies on muscular
movement and one-way valves --- a low-pressure return bus running parallel
to the main plumbing.

\section{Localized Filter Checkpoints: Lymph Nodes \texorpdfstring{(\ar{عقدة ليمفاوية})}{}}
Lymph nodes (\ar{عقدة ليمفاوية}) are the firewall's inspection
checkpoints. Lymph is forced through each node, where resident immune
cells screen it for foreign payloads. A reactive, enlarged node is a
status flag: the checkpoint is actively processing a threat. Knowing the
location of the major node clusters lets a responder read those flags on
physical exam.

\section{Healthy Operating Parameters}
The specification sheet. These are the known-good baselines against which
all Phase 3 faults are measured (feline values shown; canine ranges
differ and are noted where relevant).

\begin{center}
\renewcommand{\arraystretch}{1.35}
\begin{tabular}{l l l}
\toprule
\textbf{Parameter} & \textbf{Healthy range (feline)} & \textbf{Read-out} \\
\midrule
Core temperature   & 38.1--39.2\,\textdegree C (100.5--102.5\,\textdegree F) & Thermometer \\
Heart rate         & 140--220 beats/min & Chest / femoral pulse \\
Respiratory rate   & 20--30 breaths/min & Chest excursion \\
Capillary refill   & $<$ 2 seconds & Gum blanch test \\
Mucous membranes   & Pink, moist & Gum / inner-lip color \\
\bottomrule
\end{tabular}
\end{center}

\noindent Neonatal kittens run hotter on heart rate and are far more
vulnerable to thermal loss; a chilled kitten is a system already drifting
out of spec.

%------------------------- Chapter 3: Phase 2 ---------------------
\chapter{Phase 2 --- Pathophysiology \& System Fault Trees}

\begin{kumbox}{PHASE 2 OBJECTIVE \quad ($\approx$ 48 study hours)}
Formalize how hostile payloads execute code-injection attacks and how a
single breach cascades into total system failure. \textbf{Exit
criterion:} trace any presenting symptom backward to its originating
subsystem fault and forward to its downstream cascade.
\end{kumbox}

\section{The Threat Model: Three Payload Classes}
All acute infectious and toxic insults reduce to three payload classes,
each with a distinct attack signature.
\begin{description}[leftmargin=2em,style=nextline]
  \item[Viral payload (code injection) --- \ar{العدوى الفيروسية}]
        Hijacks the host cell's own replication machinery to manufacture
        copies of itself, destroying the host cell in the process.
  \item[Bacterial payload (resource hijack) --- \ar{العدوى البكتيرية}]
        Consumes host resources and releases endotoxins and exotoxins
        that poison surrounding tissue.
  \item[Toxic payload (corrosive input) --- \ar{السموم}]
        Inflicts direct chemical damage on cell membranes and enzyme
        systems; unlike the others, it does not replicate.
\end{description}

\section{Code-Injection Mechanics}
A virus carries no factory of its own. It is, in effect, a malicious
instruction set that must commandeer host hardware to run. It binds to a
cell, injects its genetic payload, and redirects the cell's machinery to
produce viral copies until the cell ruptures. Pathogens that depend on
host cell division therefore preferentially attack the body's
high-mitosis cores --- the same production plants catalogued in Phase 1.

\section{Immune Wipe: Feline Panleukopenia \texorpdfstring{(\ar{طاعون القطط الفيروسي})}{}}
Feline Panleukopenia (\ar{طاعون القطط الفيروسي}) is the archetypal
immune-wipe attack and the founding case of this project. The virus
targets the body's fastest-dividing cells simultaneously: the bone marrow
(\ar{نخاع العظم}), the intestinal crypt cells, and lymphoid tissue. The
result is \emph{panleukopenia} --- a collapse across all white-cell lines,
the literal meaning of the name. The firewall does not merely weaken; its
manufacturing plant goes offline.

\begin{kumbox}{FAULT TREE A --- THE IMMUNE-WIPE CASCADE}
\ttfamily\small
INGRESS: FPV (fecal-oral)\\
\quad $\downarrow$\\
TARGET: high-mitosis cores (rapidly dividing cells)\\
\quad $\downarrow$ \quad splits to three fronts:\\
\quad [1] BONE MARROW  $\rightarrow$ WBC factory OFFLINE $\rightarrow$ panleukopenia\\
\quad [2] GUT CRYPTS   $\rightarrow$ barrier BREACHED   $\rightarrow$ bacterial translocation\\
\quad [3] LYMPHOID     $\rightarrow$ firewall DEPLETED   $\rightarrow$ no antibody response\\
\quad $\downarrow$ \quad (three fronts converge)\\
SYSTEMIC SEPSIS CASCADE\\
\quad $\downarrow$\\
MULTI-ORGAN FAULT / SHOCK: liver shutdown $\cdot$ DIC $\cdot$ death
\end{kumbox}

\section{Gut-Barrier Breach \& Bacterial Translocation \texorpdfstring{(\ar{حاجز الأمعاء})}{}}
The intestinal lining is a perimeter wall (\ar{حاجز الأمعاء}) separating
the body's sterile interior from a gut lumen densely populated with
bacteria. When panleukopenia destroys the rapidly renewing cells that
maintain this wall, it is breached. Gut bacteria translocate across the
broken barrier into the bloodstream precisely when the white-cell
firewall has gone offline --- an intrusion at the exact moment the defense
grid is down.

\section{Cascade Failure: Systemic Sepsis \& Liver Shutdown}
Unchecked bacterial translocation drives \textbf{systemic sepsis}
(\ar{الإنتان الجهازي}) --- a body-wide, dysregulated inflammatory storm.
Blood pressure falls, perfusion fails, and downstream organs are starved.
The liver (\ar{الكبد}), the body's detox daemon, is overwhelmed by
endotoxin load and can shut down. Disseminated intravascular coagulation
(DIC) consumes the platelet subsystem, and the animal slides into
distributive and hypovolemic shock. This terminal cascade is what a
Phase 3 responder is racing to interrupt.

%------------------------- Chapter 4: Phase 3 ---------------------
\chapter{Phase 3 --- Field Triage \& Runtime Debugging}

\begin{kumbox}{PHASE 3 OBJECTIVE \quad ($\approx$ 30 study hours)}
Execute a deterministic diagnostic loop on an animal in acute system
failure. \textbf{Exit criterion:} run the full loop under time pressure
and produce a clean, structured handoff for the receiving veterinarian.
\textbf{Every branch converges on professional handoff.}
\end{kumbox}

\section{The Deterministic Triage Loop}
The loop is intentionally ordered: each check gates the next, mirroring an
airway-first emergency protocol. Run it top to bottom; do not jump ahead.

\begin{enumerate}[leftmargin=2.4em,label=\textbf{\arabic*.}]
  \item \textbf{[A] Airway (\ar{مجرى الهواء}).} Is the airway clear? If
        not, clear the visible obstruction and re-check.
  \item \textbf{[B] Breathing (\ar{التنفس}).} Is breathing present and
        effective? If not, support ventilation and escalate immediately.
  \item \textbf{[T] Temperature (\ar{الحرارة}).} In range? If too low
        (the common kitten failure), begin gradual external warming; if
        too high, begin controlled cooling.
  \item \textbf{[C] Circulation / CRT (\ar{زمن امتلاء الشعيرات}).} Is
        capillary refill under 2 seconds and are the gums pink? If not,
        treat for shock and poor perfusion.
  \item \textbf{[F] Fluid status (\ar{السوائل}).} Is the animal stable or
        collapsing from fluid loss? Manage collapse and conserve heat.
  \item \textbf{HANDOFF.} Transfer to a professional veterinarian with a
        structured report of every reading above.
\end{enumerate}

\section{Thermal Defense}
A collapsed kitten is most often a \emph{cold} kitten. Hypothermia slows
every downstream system and must be corrected gradually --- aggressive,
sudden re-warming is itself a fault condition. Warm the core, insulate
against further loss, and never feed a hypothermic animal whose gut has
shut down.

\section{Perfusion Debug: Capillary Refill Time}
The gum is a status LED for the circulatory plumbing. Press the mucous
membrane until it blanches white, release, and count the seconds to
re-perfusion.

\begin{center}
\renewcommand{\arraystretch}{1.35}
\begin{tabular}{l l}
\toprule
\textbf{Gum / CRT read-out} & \textbf{System interpretation} \\
\midrule
Pink, refill $<$ 2 s        & Perfusion nominal \\
Pale / white, refill $>$ 2 s & Poor perfusion, shock, or anemia \\
Blue (cyanotic)             & Oxygenation failure --- escalate now \\
Brick red                   & Hyperdynamic sepsis \\
Yellow (icteric)            & Liver / red-cell breakdown \\
\bottomrule
\end{tabular}
\end{center}

\section{Fluid-Collapse Management: The Deterministic Math}
Fluid loss is quantifiable. A responder estimates the dehydration
percentage from clinical signs, then computes the fluid deficit so the
receiving veterinarian inherits a number, not a guess. Let $BW$ be body
weight in kilograms and $D$ the estimated dehydration as a percentage:

\begin{equation}
  V_{\text{deficit}}\;[\text{mL}] \;=\; BW\;[\text{kg}] \;\times\; D\;[\%] \;\times\; 10 .
\end{equation}

\noindent The shock index gives a fast, single-number read on
cardiovascular stability, where $\mathrm{HR}$ is heart rate and
$\mathrm{SAP}$ is systolic arterial pressure:

\begin{equation}
  \mathrm{SI} \;=\; \frac{\mathrm{HR}}{\mathrm{SAP}} .
\end{equation}

\noindent A rising shock index signals a circulatory system trending
toward collapse. Both figures belong in the handoff report.

\begin{kumbox}{HANDOFF CONTRACT}
The triage loop does not cure --- it buys time and produces data. The
deliverable at the end of every run is a structured handoff: the
known-good baselines, the observed deviations, the interventions applied,
and the computed deficit. \textbf{Stabilize, document, hand off.}
\end{kumbox}

%------------------------- Back matter ----------------------------
\backmatter
\chapter{Closing Note}
Project Kumquat treats biology as an engineering problem not to diminish
its complexity, but to make that complexity survivable in the field. The
framework is open and bilingual by design so that it can be taught,
forked, and translated. It honors a tiny fighter, Koki, by ensuring the
next responder has a loop to run instead of a guess to make.

\begin{center}
\vspace{1em}
{\itshape\color{kumdeep}When a system goes down, we do not guess --- we debug.}\\[0.4em]
{\sffamily\color{kumgrey} In memory of Kumquat (Koki) \quad \ar{طاعون القطط الفيروسي}}
\end{center}

\end{document}

%==================================================================
%  APPENDIX: LEGACY FALLBACK (pdfLaTeX + arabtex)
%------------------------------------------------------------------
%  If a Unicode engine is unavailable and you must use pdfLaTeX, the
%  arabtex package can render Arabic from ASCII transliteration. Replace
%  the multilingual block above with:
%
%     \usepackage[T1]{fontenc}
%     \usepackage{arabtex}
%     \usepackage{utf8}            % arabtex's own utf8 input layer
%     \setcode{utf8}
%     % inline usage:  \RL{\<...arabic transliteration...>}
%
%  NOTE: arabtex expects either transliterated input or its own utf8
%  layer and does not consume native Unicode as cleanly as polyglossia.
%  For "flawless inline Arabic," XeLaTeX/LuaLaTeX + polyglossia (the
%  default configuration above) is strongly recommended.
%==================================================================
